% Generated draft: requires mathematical and source review.
\documentclass[11pt]{article}
\usepackage[margin=25mm]{geometry}
\usepackage{fontspec}
\setmainfont{Latin Modern Roman}
\usepackage{amsmath,amssymb,mathtools,unicode-math}
\setmathfont{Latin Modern Math}
\usepackage{enumitem,needspace,xurl,hyperref}
\setlength{\parindent}{0pt}
\setlength{\parskip}{5pt}
\setlength{\emergencystretch}{3em}
\newcommand{\N}{\mathbb{N}}
\newcommand{\Z}{\mathbb{Z}}
\newcommand{\Q}{\mathbb{Q}}
\newcommand{\R}{\mathbb{R}}
\newcommand{\C}{\mathbb{C}}
\newcommand{\F}{\mathbb{F}}
\newcommand{\A}{\mathbb{A}}
\newcommand{\PP}{\mathbb{P}}
\newcommand{\E}{\mathbb{E}}
\newcommand{\eps}{\varepsilon}
\newcommand{\dd}{\,\mathrm{d}}
\begin{document}
\section*{UnsolvedMath ID: 20003325}
\textbf{Source problem code:} catalogue entry
\subsection{Definitions and mathematical statement}
The source's Conjecture 38 (quantum unique ergodicity) states that negative sectional curvature of \(X\) should imply \(\mu_i\to dV\). The excerpt does not supply the preceding definitions of \(\mu_i\), the normalization of \(dV\), compactness or finite-volume hypotheses, or the eigenvalue enumeration. A standard configuration-space interpretation, conditional on these missing conventions, is the following. Let \(X\) be a complete Riemannian manifold of finite positive volume, put \(dV=d\operatorname{vol}_X/\operatorname{vol}(X)\), and let \(\varphi_i\) be normalized \(L^2(X,dV)\) Laplace eigenfunctions with eigenvalues tending to infinity. Define \(d\mu_i=|\varphi_i|^2dV\). Under the source's negative-sectional-curvature hypothesis, the conjectured conclusion is \(\int_X a\,d\mu_i\to\int_X a\,dV\) for every bounded continuous test function \(a\). On a noncompact manifold this formulation includes absence of escape of mass. The source does not determine whether its intended class is compact, finite-volume noncompact, or otherwise; an infinite-volume space has no normalized Riemannian probability measure of this form. No assertion about continuous-spectrum generalized eigenfunctions or microlocal lifts is specified here.\par The arithmetic specialization in the source takes \(X=\Gamma\backslash\widetilde X\), with \(\widetilde X\) a symmetric space of noncompact type and \(\Gamma\) an arithmetic lattice, and requires joint eigenfunctions of the invariant differential operators and Hecke operators. It calls weak-* limits of the associated \(\mu_i\) arithmetic quantum limits. The source describes the Laplacian as the relevant invariant differential operator in rank one and cites [32] for the belief that \(\mathrm{SL}_2(\Z)\backslash\mathbb H^2\) has simple cuspidal Laplace spectrum, which would make the Hecke eigenfunction condition automatic on those eigenspaces. This belief is a supplied source qualification, not a verified present-day status claim. The broader symmetric-space discussion is not itself a claim that every such space has strictly negative sectional curvature.\par The supplied research aid proposes a conditional, basis-independent certificate. For a finite-dimensional Laplace eigenspace \(E_\lambda\subset L^2(X,dV)\), let \(m_\lambda=\dim E_\lambda\), let \(P_\lambda\) be its orthogonal projection, and take a real bounded observable \(a\). Write \(\bar a=\int_X a\,dV\), \(M_b f=bf\), and\[T_\lambda(a)=P_\lambda M_{a-\bar a}P_\lambda\big|_{E_\lambda},\qquad D_\lambda(a)=\sup_{\substack{\varphi\in E_\lambda\\\|\varphi\|_2=1}}\left|\int_X a|\varphi|^2dV-\bar a\right|,\qquad V_\lambda(a)=\frac{\operatorname{tr}(T_\lambda(a)^2)}{m_\lambda}.\] Its claimed identities and sharp multiplicity bounds are\[D_\lambda(a)=\|T_\lambda(a)\|_{\mathrm{op}},\qquad \sqrt{V_\lambda(a)}\le D_\lambda(a)\le\sqrt{m_\lambda V_\lambda(a)}.\] Thus \(m_\lambda V_\lambda(a)\to0\) on a countable uniformly dense family of real continuous observables is a sufficient all-unit-eigenvector QUE criterion, with the compact case using \(C(X,\R)\) and the locally compact, second-countable noncompact case using \(C_0(X,\R)\). The latter uses the fact that probability measures converging on a uniformly dense subset of \(C_0(X)\) to the integrals against a specified probability measure converge weakly to that probability measure and have no escape of mass. A local limit proportional to volume with an unspecified coefficient does not provide that conclusion. The variance here includes all matrix entries: vanishing of only the diagonal variance in one chosen basis does not control \(D_\lambda\), as the abstract self-adjoint matrix \(\left(\begin{smallmatrix}0&1\\1&0\end{smallmatrix}\right)\) illustrates. This matrix illustration is not a constructed counterexample to QUE. These conditional certificates do not establish the conjecture's required estimates for negatively curved manifolds.
\subsection{Short English statement}
The conjecture asks for equidistribution of every high-energy eigenfunction's probability density. The supplied aid gives a criterion that controls all vectors in a multiple eigenspace and handles loss of mass in a cusp. Missing normalization and geometric hypotheses remain explicit; neither the criterion nor the recorded progress labels constitute a solution of QUE.
\subsection{Sources}
\begin{itemize}\item Local UnsolvedMath registry, permanent ID 20003325, code catalogue entry; canonical file \texttt{aim-other-notes.json}, \texttt{notes[131]}, conjecture tag.\item American Institute of Mathematics Problem Lists, Other domain, workshop \emph{Emerging applications of measure rigidity}, source item 38, Conjecture 38 (Quantum unique ergodicity); no section supplied. \url{https://aimath.org/WWN/measrigid/measrigid.pdf}.\item Source reference [32], cited for the conjectured simplicity of the cuspidal spectrum of \(\mathrm{SL}_2(\Z)\backslash\mathbb H^2\); its bibliographic entry is absent from this record.\item Attached research attempt 1, titled \emph{A multiplicity- and cusp-safe QUE certificate}, recorded 2026-08-14: result \texttt{partial\_result}, problem-at-run \texttt{open}, rigor \texttt{proved}, novelty confidence low; registry status \texttt{partially\_solved}, classification \texttt{PARTIAL-PROGRESS}. This machine-generated research aid is not peer reviewed and requires independent verification before citation.\end{itemize}
\end{document}
